\documentclass[a4paper,12pt]{article}
\usepackage[T2A]{fontenc}
\usepackage[utf8]{inputenc}
\usepackage[russian]{babel}
\usepackage{amsmath,amssymb,mathtools}
\usepackage{helvet}
\renewcommand{\familydefault}{\sfdefault}
\usepackage{icomma}
\usepackage[a4paper,left=18mm,right=18mm,top=18mm,bottom=20mm]{geometry}
\usepackage{microtype}
\usepackage{tikz}
\usepackage{pgfplots}
\pgfplotsset{compat=1.18}
\usepackage{tabularx,booktabs,longtable,array}
\usepackage{enumitem}
\usepackage{hyperref}
\usepackage{parskip}
\usepackage{fancyhdr}
\usepackage{setspace}
\usepackage{xcolor}

\definecolor{accent}{HTML}{287F7A}
\definecolor{darkaccent}{HTML}{1C3151}
\definecolor{textgray}{HTML}{333333}
\definecolor{softgray}{HTML}{667184}
\definecolor{coral}{HTML}{C8644B}

\hypersetup{colorlinks=true,linkcolor=darkaccent,urlcolor=accent}
\onehalfspacing
\setlength{\parindent}{0pt}
\setlength{\emergencystretch}{2em}
\setlength{\headheight}{15pt}
\setlist[itemize]{leftmargin=*,itemsep=3pt,topsep=4pt}
\setlist[enumerate]{leftmargin=*,itemsep=6pt,topsep=4pt}
\pagestyle{fancy}
\fancyhf{}
\lhead{\textcolor{textgray}{Линейная функция}}
\rhead{\textcolor{textgray}{София Кальней}}
\cfoot{\textcolor{softgray}{\thepage}}
\renewcommand{\headrulewidth}{0.4pt}

\newcommand{\sectionline}{\par\vspace{2pt}\textcolor{accent}{\rule{\linewidth}{0.8pt}}\vspace{5pt}}
\newcommand{\key}[1]{\textcolor{darkaccent}{\textbf{#1}}}
\newcommand{\warn}[1]{\textcolor{coral}{\textbf{#1}}}

\begin{document}

\begin{titlepage}
  \centering
  \vspace*{24mm}
  {\color{accent}\Large ЧЕК-ЛИСТ\par}
  \vspace{8mm}
  {\color{darkaccent}\bfseries\fontsize{30}{36}\selectfont
  График линейной функции\par}
  \vspace{5mm}
  {\Large Формула, построение и уравнение прямой\par}
  \vfill
  {\large Лёвин Артём Александрович\par}
  \vspace{2mm}
  {\large эксклюзивно для Софьи Кальней\par}
  \vspace{8mm}
  {\large \today\par}
  \vfill
  \begin{tikzpicture}[x=1cm,y=1cm]
    \draw[accent,line width=1.3pt]
      (-6,0) .. controls (-3.8,1.2) and (-1.7,-1.0) .. (0,0.15)
      .. controls (2.2,1.25) and (4.0,-0.8) .. (6,0.35);
    \fill[coral] (-6,0) circle (2.2pt);
    \fill[coral] (6,0.35) circle (2.2pt);
  \end{tikzpicture}
\end{titlepage}

\section*{1. Главная модель}
\sectionline

\[
  \boxed{y=kx+b}
\]

\begin{itemize}
  \item \key{Линейная функция} задаётся формулой $y=kx+b$, где $k$ и $b$ --- числа.
  \item Её график --- \key{прямая}. Для построения прямой достаточно двух различных точек.
  \item $k$ --- \key{угловой коэффициент}: он отвечает за направление и крутизну наклона.
  \item $b$ --- \key{ордината пересечения с осью $Oy$}: при $x=0$ получаем точку $(0;b)$.
\end{itemize}

\begin{center}
\begin{tikzpicture}
\begin{axis}[
  width=11cm,height=7cm,axis lines=middle,unit vector ratio=1 1,
  xmin=-4,xmax=5,ymin=-4,ymax=6,samples=200,clip=false,
  xlabel={$x$},ylabel={$y$},grid=both,minor tick num=1,
  grid style={draw=softgray!18},tick label style={font=\small}
]
  \addplot[accent,very thick,domain=-4:5]{1.2*x+1};
  \addplot[only marks,mark=*,coral] coordinates {(0,1)};
  \node[anchor=west,text=coral] at (axis cs:0.15,1.2) {$(0;b)$};
  \node[anchor=south east,text=accent] at (axis cs:4.2,5.7) {$y=1{,}2x+1$};
\end{axis}
\end{tikzpicture}
\end{center}

\section*{2. Как коэффициент $k$ меняет прямую}
\sectionline

\begin{center}
\renewcommand{\arraystretch}{1.4}
\begin{tabularx}{\linewidth}{@{}>{\bfseries}l X X@{}}
\toprule
Условие & Поведение функции & Что видно на графике \\
\midrule
$k>0$ & функция возрастает & прямая поднимается слева направо \\
$k<0$ & функция убывает & прямая опускается слева направо \\
$k=0$ & $y=b$ --- постоянная функция & горизонтальная прямая \\
\bottomrule
\end{tabularx}
\end{center}

Если известны две точки $A(x_1;y_1)$ и $B(x_2;y_2)$, то
\[
  \boxed{k=\frac{y_2-y_1}{x_2-x_1}},\qquad x_2\ne x_1.
\]

Числитель показывает изменение $y$, знаменатель --- изменение $x$. Условие
$x_2\ne x_1$ обязательно: вертикальная прямая не является графиком функции
вида $y=kx+b$.

\section*{3. Алгоритм построения по формуле}
\sectionline

\begin{enumerate}
  \item Выберите два удобных различных значения $x$.
  \item Подставьте каждое значение в формулу и вычислите $y$.
  \item Запишите пары координат $(x;y)$ в таблицу.
  \item Отметьте обе точки на координатной плоскости.
  \item Проведите через них единственную прямую и подпишите формулу.
\end{enumerate}

\key{Пример 1.} Построим $y=-2x+3$.

\begin{center}
\renewcommand{\arraystretch}{1.35}
\begin{tabular}{@{}ccc@{}}
\toprule
$x$ & $0$ & $2$ \\
\midrule
$y=-2x+3$ & $3$ & $-1$ \\
\bottomrule
\end{tabular}
\end{center}

Получены точки $(0;3)$ и $(2;-1)$. Проверка: $k=-2<0$, поэтому прямая
действительно должна убывать.

\begin{center}
\begin{tikzpicture}
\begin{axis}[
  width=10cm,height=7cm,axis lines=middle,unit vector ratio=1 1,
  xmin=-2,xmax=4,ymin=-4,ymax=6,samples=200,clip=false,
  xlabel={$x$},ylabel={$y$},grid=both,grid style={draw=softgray!18}
]
  \addplot[coral,very thick,domain=-2:4]{-2*x+3};
  \addplot[only marks,mark=*,accent] coordinates {(0,3) (2,-1)};
  \node[anchor=south west] at (axis cs:0,3) {$(0;3)$};
  \node[anchor=north west] at (axis cs:2,-1) {$(2;-1)$};
\end{axis}
\end{tikzpicture}
\end{center}

\section*{4. Как восстановить уравнение прямой}
\sectionline

\subsection*{4.1. По двум точкам}

Для точек $A(-1;5)$ и $B(3;-3)$:
\[
  k=\frac{-3-5}{3-(-1)}=\frac{-8}{4}=-2.
\]
Подставим $A$ в $y=kx+b$:
\[
  5=-2\cdot(-1)+b,\qquad b=3.
\]
Итак, \key{$y=-2x+3$}. Подстановка координат точки $B$ даёт
$-3=-6+3$, значит уравнение найдено верно.

\subsection*{4.2. По угловому коэффициенту и одной точке}

Если прямая имеет наклон $k$ и проходит через $(x_0;y_0)$, удобно использовать
формулу
\[
  \boxed{y-y_0=k(x-x_0)}
  \qquad\Longleftrightarrow\qquad
  \boxed{y=k(x-x_0)+y_0}.
\]

\key{Пример 2.} Прямая с $k=3$ проходит через $R(2;-1)$:
\[
  y=3(x-2)-1=3x-7.
\]

\section*{5. Взаимное положение и параметр}
\sectionline

\begin{itemize}
  \item Прямые $y=k_1x+b_1$ и $y=k_2x+b_2$ \key{параллельны}, если
  $k_1=k_2$ и $b_1\ne b_2$.
  \item Прямые \key{совпадают}, если $k_1=k_2$ и $b_1=b_2$.
  \item При $k_1\ne k_2$ пересечение единственно. Для его поиска приравнивают
  правые части, находят $x$, затем вычисляют $y$.
\end{itemize}

\key{Пример 3.} Найдём пересечение $y=2x-1$ и $y=-x+5$:
\[
  \begin{aligned}
    2x-1&=-x+5, & 3x&=6, & x&=2,\\
    y&=2\cdot2-1=3.
  \end{aligned}
\]
Точка пересечения: $(2;3)$.

В семействе $y=2x+a$ угловой коэффициент остаётся равным $2$, поэтому все
прямые параллельны. Параметр $a$ задаёт точку $(0;a)$ и вертикальный сдвиг
графика.

\section*{6. Типичные ошибки}
\sectionline

\begin{enumerate}
  \item \warn{Путать $k$ и $b$.} Сначала найдите точку $(0;b)$, затем учитывайте наклон $k$.
  \item \warn{Менять порядок вычитания только в числителе.} В формуле для $k$ порядок точек должен быть одинаковым сверху и снизу.
  \item \warn{Строить прямую по одной точке.} Через одну точку проходит бесконечно много прямых.
  \item \warn{Терять скобки у отрицательных координат.} Например, $3-(-1)=4$.
  \item \warn{Не проверять ответ второй точкой.} Подстановка быстро обнаруживает ошибку в $k$ или $b$.
\end{enumerate}

\section*{7. Итоговый чек-лист}
\sectionline

\begingroup
\setstretch{1.22}
\begin{itemize}[label=$\square$,itemsep=2pt]
  \item Я узнаю линейную функцию по формуле $y=kx+b$.
  \item Я объясняю смысл коэффициентов $k$ и $b$.
  \item Я определяю возрастание или убывание по знаку $k$.
  \item Я строю прямую по двум точкам и проверяю вычисления.
  \item Я нахожу $k$ по координатам двух точек.
  \item Я составляю уравнение через точку при известном $k$.
  \item Я нахожу точку пересечения двух прямых.
  \item Я распознаю параллельные и совпадающие прямые.
  \item Я понимаю вертикальный сдвиг в семействе $y=kx+a$.
\end{itemize}
\endgroup

\clearpage
\section*{Тренировочный блок}
\sectionline

\textcolor{softgray}{Сначала решите задания самостоятельно. Ответы находятся на следующей странице.}

\subsection*{Средний уровень}
\begin{enumerate}
  \item Для функции $y=2x-3$ заполните таблицу при $x=-1$, $0$, $2$.
  \item Найдите точки пересечения графика $y=-3x+6$ с осями координат.
  \item Определите, возрастает или убывает функция $y=-\frac12x+4$, и укажите точку пересечения с осью $Oy$.
\end{enumerate}

\subsection*{Уровень выше среднего}
\begin{enumerate}[resume]
  \item Составьте уравнение прямой, проходящей через $A(-2;5)$ и $B(4;-1)$.
  \item Составьте уравнение прямой с угловым коэффициентом $4$, проходящей через точку $(3;-2)$.
  \item Найдите точку пересечения прямых $y=2x-5$ и $y=-x+7$.
  \item Найдите $a$, если прямая $y=3x+a$ проходит через точку $(-2;5)$.
  \item Найдите $a$, если прямая $y=(a-1)x+4$ параллельна прямой $y=5x-2$.
  \item Прямые $y=-2x+1$ и $y=x-5$ пересекаются в точке $R$. Составьте уравнение прямой, проходящей через $R$ параллельно $y=4x+7$.
\end{enumerate}

\subsection*{Сложный уровень}
\begin{enumerate}[resume]
  \item Найдите $a$, если точка пересечения прямых
  $y=(a-1)x+2a$ и $y=3x-4$ имеет абсциссу $2$. Запишите обе прямые при найденном $a$ и проверьте точку пересечения.
\end{enumerate}

\clearpage
\section*{Ответы}
\sectionline

\renewcommand{\arraystretch}{1.45}
\begin{longtable}{@{}>{\bfseries}p{12mm}p{\dimexpr\linewidth-18mm\relax}@{}}
\toprule
№ & Ответ \\
\midrule
\endhead
1 & $(-1;-5)$, $(0;-3)$, $(2;1)$ \\
2 & $(0;6)$ и $(2;0)$ \\
3 & Убывает; $(0;4)$ \\
4 & $y=-x+3$ \\
5 & $y=4x-14$ \\
6 & $(4;3)$ \\
7 & $a=11$ \\
8 & $a=6$ \\
9 & $R(2;-3)$; $y=4x-11$ \\
10 & $a=1$; $y=2$ и $y=3x-4$; точка $(2;2)$ \\
\bottomrule
\end{longtable}

\vfill
\begin{center}
  \textcolor{softgray}{Код самопроверки: \textbf{2416}}
\end{center}

\end{document}
